% Section 4.4: Blades | Part: 1 | Depends on: sec:tensors, sec:metric | Status: done | Original lines: 1113-1175
%!TEX root = ../../main.tex
\subsection{Blades}\label{sec:blades}
\begin{definition}[Blade]
A $k$-vector of the form $A=a_1\wedge\cdots\wedge a_k$ is called a \emph{$k$-blade} (or \emph{simple}, \emph{decomposable} $k$-vector).
A general $k$-vector is a sum of $k$-blades.
\end{definition}

\begin{proposition}[Blades $\leftrightarrow$ subspaces]\label{prop:blade}
Let $A=a_1\wedge\cdots\wedge a_k\neq0$. Then $U_A:=\{v\in V: v\wedge A=0\}=\Span(a_1,\dots,a_k)$, a $k$-dimensional subspace.
Two nonzero blades define the same subspace iff they are proportional.
\end{proposition}
\begin{proof}
By Theorem~\ref{thm:extbasis}, $v\wedge a_1\wedge\cdots\wedge a_k=0$ iff $v,a_1,\dots,a_k$ are linearly dependent, iff $v\in\Span(a_i)$. If $b_1\wedge\cdots\wedge b_k$ spans the same subspace, $b_j=M_j^{\,i}a_i$ with $M$ invertible, and Corollary~\ref{cor:det} gives $b_1\wedge\cdots\wedge b_k=\det M\;a_1\wedge\cdots\wedge a_k$.
\end{proof}

Not every $k$-vector is a blade. For instance, $e_1\wedge e_2+e_3\wedge e_4\in\bigwedge^2\R^4$ is not (it satisfies $A\wedge A=2\,e_1\wedge e_2\wedge e_3\wedge e_4\neq0$, whereas $a\wedge b\wedge a\wedge b=0$ for any blade). In fact a nonzero $2$-vector is a blade iff $A\wedge A=0$ (Pl\"ucker relation).

% Figure | used in parts/01-algebra/04-exterior-algebra/04-blades.tex
%!TEX root = ../main.tex
\begin{figure}[ht]
\centering
\begin{tikzpicture}[scale=1.15,>=Stealth]
  % left: a wedge b
  \coordinate (O) at (0,0);
  \coordinate (a) at (2.6,0);
  \coordinate (b) at (1.0,1.7);
  \fill[blue!15] (O)--(a)--($(a)+(b)$)--(b)--cycle;
  \draw[thick,->] (O)--(a) node[below]{$a$};
  \draw[thick,->] (O)--(b) node[left]{$b$};
  \draw[dashed] (a)--($(a)+(b)$)--(b);
  \draw[->,red,thick] (1.35,0.55) arc[start angle=-30,end angle=210,radius=0.38];
  \node at (1.8,-0.85) {$a\wedge b$};
  % right: sheared
  \begin{scope}[xshift=6.3cm]
    \coordinate (O) at (0,0);
    \coordinate (a) at (2.6,0);
    \coordinate (b2) at (2.6,1.7); % b + a*(1.6/2.6)... shear
    \fill[blue!15] (O)--(a)--($(a)+(b2)$)--(b2)--cycle;
    \draw[thick,->] (O)--(a) node[below]{$a$};
    \draw[thick,->] (O)--(b2) node[above left]{$b'=b+\lambda a$};
    \draw[dashed] (a)--($(a)+(b2)$)--(b2);
    \draw[->,red,thick] (1.35,0.55) arc[start angle=-30,end angle=210,radius=0.38];
    \node at (2.6,-0.85) {$a\wedge b'=a\wedge b$};
  \end{scope}
\end{tikzpicture}
\caption{The bivector $a\wedge b$ is an oriented area. Shearing $b\mapsto b+\lambda a$ changes the parallelogram but not the bivector, since $a\wedge a=0$.}
\label{fig:wedge}
\end{figure}


\begin{intuition}
\begin{itemize}
  \item A \textbf{vector} is a directed length (an oriented line segment). A \textbf{bivector} $a\wedge b$ is a \emph{directed area}: an
  oriented patch of plane with a magnitude (the area of the parallelogram) and a sense of circulation (go along $a$, then along $b$);
  see Fig.~\ref{fig:wedge}. A \textbf{trivector} $a\wedge b\wedge c$ is an oriented volume, and so on: grade $0,1,2,3,\dots$ = point (scalar), line, plane, volume.
  \item $a\wedge b=0$ exactly when $a,b$ are parallel: the parallelogram is degenerate. This is the algebraic statement that ``zero area'' means ``dependent''.
  \item Only the \emph{plane}, \emph{area} and \emph{orientation} matter, not the particular shape of the parallelogram: $a\wedge(b+\lambda a)=a\wedge b$.
  \item $a\wedge b=-b\wedge a$: reversing the circulation reverses the orientation, just like swapping two columns flips the sign of a determinant (Corollary~\ref{cor:det}: the determinant \emph{is} the volume of the top-grade wedge).
  \item Unlike the cross product, which only exists in $3$D and secretly produces a vector perpendicular to a plane, the wedge works in every dimension and keeps the plane \emph{as a plane}. (We recover the cross product through duality in \S\ref{sec:duality}.)
  \item The exterior algebra knows nothing about lengths or angles: it is metric-free. To also measure, we need the Clifford algebra.
\end{itemize}
\end{intuition}
